% Generated draft: requires mathematical and source review.
\documentclass[11pt]{article}
\usepackage[margin=25mm]{geometry}
\usepackage{fontspec}
\setmainfont{Latin Modern Roman}
\usepackage{amsmath,amssymb,mathtools,unicode-math}
\setmathfont{Latin Modern Math}
\usepackage{enumitem,needspace,xurl,hyperref}
\setlength{\parindent}{0pt}
\setlength{\parskip}{5pt}
\setlength{\emergencystretch}{3em}
\newcommand{\N}{\mathbb{N}}
\newcommand{\Z}{\mathbb{Z}}
\newcommand{\Q}{\mathbb{Q}}
\newcommand{\R}{\mathbb{R}}
\newcommand{\C}{\mathbb{C}}
\newcommand{\F}{\mathbb{F}}
\newcommand{\A}{\mathbb{A}}
\newcommand{\PP}{\mathbb{P}}
\newcommand{\E}{\mathbb{E}}
\newcommand{\eps}{\varepsilon}
\newcommand{\dd}{\,\mathrm{d}}
\begin{document}
\section*{UnsolvedMath ID: 30004319}
\textbf{Source problem code:} catalogue entry
\subsection{Definitions and mathematical statement}
For a possibly nonassociative ring \(R\), write \((x,y,z)=(xy)z-x(yz)\) for its associator. The ring is \emph{alternative} if it satisfies both alternative identities
\[
(xx)y=x(xy),\qquad y(xx)=(yx)x
\qquad\text{for every }x,y\in R.
\]
Equivalently, \((x,x,y)=(y,x,x)=0\) for all \(x,y\in R\). No restriction on characteristic is needed to state these identities.

\textbf{Question as supplied.} Is every ring \(R\) that parametrizes an \(A_2\)-graded group alternative? Formally, let \(\operatorname{Par}_{A_2}(R,G)\) denote the source's intended assertion that \(R\) parametrizes the \(A_2\)-graded group \(G\). The requested assertion is
\[
\forall R\;\Bigl[
 \bigl(\exists G\;\operatorname{Par}_{A_2}(R,G)\bigr)
 \Longrightarrow
 \bigl(\forall x,y\in R:\ (xx)y=x(xy)\ \text{and}\ y(xx)=(yx)x\bigr)
\Bigr].
\]
Here \(R\) ranges over the rings permitted by the intended parametrization framework. The notation \(\operatorname{Par}_{A_2}\) is only a placeholder: the supplied record gives neither the axioms of an \(A_2\)-graded group nor the definition of parametrization. It also does not specify whether the rings must have an identity or satisfy additional hypotheses. Consequently, the displayed assertion records the intended quantifiers and conclusion but is not a fully specified standalone formulation. Associativity must not be silently imposed, since associative rings already satisfy the displayed identities.
\subsection{Short English statement}
The question asks whether every ring used to parametrize an \(A_2\)-graded group satisfies the left and right alternative identities. The source's grading and parametrization conventions are missing, so the precise class of rings covered by the question requires human review.
\subsection{Sources}
\begin{itemize}
\item Local UnsolvedMath registry: supplied record with permanent ID \(30004319\), title \emph{Alternativity of Parameters for \(A_2\)-Graded Groups}, exact source problem code \texttt{catalogue entry}, in the Oberwolfach Reports Open Problems collection. No individual proposer is identified. The supplied proposed year is 2019.
\item Supplied report citation: \emph{Mini-Workshop: Rank One Groups and Exceptional Algebraic Groups} (2020), Oberwolfach Reports, pp.~3243--3269, DOI \texttt{10.4171/owr/2019/52}. \url{https://doi.org/10.4171/owr/2019/52}. The background describes these reports as documenting mathematical workshops and their open-problem discussions.
\item Supplied literature source: Wiedemann, \emph{Root Graded Groups}. \url{https://arxiv.org/abs/2404.02042}.
\item Supplied literature source: \emph{Root graded groups revisited}. \url{https://arxiv.org/abs/2406.03558}.
\item Supplied literature source: \emph{Oberwolfach Reports, \(A_2\)-graded groups}. \url{https://ems.press/content/serial-article-files/46832}.
\item Supplied dated curation, checked at \texttt{2026-08-22}: the record labels the problem open and states that root-graded group theory provides partial structural results, while no nonalternative parametrizing-ring example and no proof of universal alternativity are known. This is supplied dated literature triage derived from the OWR curation, not an independently verified assessment; its time-sensitive status and cited claims require independent verification.
\item The registry supplies the statement-status label \texttt{corrected\_verified}, with an explanation that mathematical content was preserved while notation, typography, omitted definitions, or duplicate clauses were repaired for self-contained publication. This supplied editorial label does not establish mathematical correctness or resolve the missing definitions in this draft.
\end{itemize}
\end{document}
